\documentclass[12pt]{article}
\usepackage[margin=1in]{geometry}
\usepackage{enumitem}
\usepackage{hyperref}

\pagestyle{empty}
\begin{document}

\begin{flushright}
Huang Zhongchang\\
Independent Researcher\\
Nanjing, China\\
\today
\end{flushright}

\vspace{0.3in}

\noindent
The Editor\\
Physical Review D\\
American Physical Society

\vspace{0.2in}

\noindent
\textbf{Re: Submission to Physical Review D}

\vspace{0.1in}

\noindent
Dear Editor,

\vspace{0.1in}

I am submitting the manuscript entitled ``Two Rules, One Universe: Gravity and Quantum Decoherence from Information Capacity'' for consideration for publication in Physical Review D.

The paper proposes that both Newtonian gravity and quantum decoherence follow from a single mechanism: the bounded information capacity of physical cells. Only two physical axioms enter the theory---causality (A1) and a 1-bit capacity bound per cell (A2)---plus one structural choice for the dynamical variable. From these, we derive the unified telegraph equation
$\partial_t^2 q + \gamma_0(1-q)\partial_t q = c^2\nabla^2(\ln q)$
for the free-capacity fraction $q$. The same equation governs all regimes without patching: the static limit yields $\nabla^2(\ln 1/q)=0$, whose spherical solution $q(r)=e^{-GM/rc^2}$ recovers Newtonian gravity; the $\gamma\to0$ vacuum limit dynamically selects Lorentz invariance as the ground state of the empty universe rather than imposing it as an axiom.

The theory makes two falsifiable predictions: a wave--diffusion crossover in gravitational dynamics at $k=\gamma/(2c)$, and a soft boundary replacing the classical event horizon. An H-theorem is proved for the full nonlinear dynamics via a Lyapunov functional, establishing thermodynamic consistency.

A distinctive feature of the work is its treatment of symmetry emergence. Where Bassani and Magueijo (2025, PRD 111, 103529) achieve diffeomorphism invariance through a Markov chain with a preferred foliation, local diffeomorphism invariance imposed pre-emergence, and five Tier-2 structural choices, the present framework selects Lorentz invariance deterministically from two axioms and three structural choices. The comparison uses the same tiered-assumption methodology, applied impartially to both theories. The economy of assumptions---fewer axioms, deterministic rather than stochastic emergence---is a central methodological point of the paper.

The quantum-mechanical foundation builds on Zilly (2026, arXiv:2603.11900), who proved that finite information capacity plus self-referential consistency yields standard quantum mechanics. We accept Zilly's theorem as the QM foundation and build upward toward gravity, so the QM-gravity bridge is constructed from a shared information-theoretic root rather than postulating a unification.

The paper is honest about its limitations. It recovers only the Newtonian limit of GR, not general relativity; the 1-bit-per-cell bound implies $S\sim R^3$ volume scaling rather than holographic area scaling; $\gamma(q)=\gamma_0(1-q)$ is a modeling ansatz rather than a rigorous derivation; and spatial dimensionality $d=3$ enters as an empirical input rather than a prediction. These are flagged prominently and not buried in fine print.

I believe Physical Review D is the appropriate venue because the paper engages directly with PRD's readership: the comparison with Bassani \& Magueijo (PRD 2025) extends a conversation already taking place in this journal, and the paper's methodology---constructing gravitational phenomenology from information-theoretic first principles---addresses foundational questions of the kind PRD regularly publishes.

I suggest the following potential referees with relevant expertise:
\begin{itemize}[nosep,leftmargin=*]
    \item Researchers in emergent gravity and entropic gravity (e.g., those familiar with Jacobson 1995, Verlinde 2011, and Padmanabhan 2010)
    \item Researchers in quantum foundations, particularly recent work on deriving QM from information-theoretic axioms
    \item Researchers in causal set theory and discrete spacetime approaches
\end{itemize}

This manuscript is not under consideration elsewhere. There are no competing interests to declare.

\vspace{0.15in}

\noindent
Respectfully submitted,

\vspace{0.3in}

\noindent
Huang Zhongchang

\end{document}
